%! Composition of Functions — Unit 1, Lesson 1.5 — Guided Notes
\documentclass[10pt]{article}
\usepackage{precalculus-article}
\usepackage{precalculus-boxes}
\newcommand{\TallMath}[1]{$\displaystyle #1\rule[-1.4em]{0pt}{3.2em}$}

\begin{document}

\pageheader{Unit 1, Lesson 1.5}{Guided Notes}

\vspace{0.12in}

\begin{objectivebox}
By the end of this lesson, I will be able to\ldots
\begin{enumerate}[itemsep=2pt]
  \item Evaluate $f(g(x))$ at a number by working from the \blank{1.8cm} out
        --- the \blank{1.6cm} function runs first.
  \item Build $(f \circ g)(x)$ by \blank{2.6cm} $g(x)$ for every $x$ in
        $f(x)$, and explain why $f(g(x))$ and $g(f(x))$ are usually
        \blank{2.2cm}.
  \item \blank{2.6cm} a function into an inner and an outer piece, and
        recognize a horizontal shift as a \blank{2.6cm}.
\end{enumerate}
\end{objectivebox}

\vspace{0.1in}

\boxguard
\begin{vocabbox}
Fill in each term as we define it together.
\par\vspace{2pt}
\noindent\textbf{\textcolor{plum}{Composition:}}\\[1pt]
\writeline\\[3pt]
\noindent\textbf{\textcolor{plum}{Inner and outer function:}}\\[1pt]
\writeline\\[3pt]
\noindent\textbf{\textcolor{plum}{Identity function:}}\\[1pt]
\writeline\\[3pt]
\noindent\textbf{\textcolor{plum}{Decomposition:}}\\[1pt]
\writeline\\[3pt]
\noindent\textbf{\textcolor{plum}{Domain of a composite:}}\\[1pt]
\writeline
\end{vocabbox}

\vspace{0.1in}

\boxguard
\begin{hookbox}
A \$60 jacket, a \$10 coupon, and 6\% sales tax. Your clerk takes the coupon
off first, then adds the tax. The clerk at the next register adds the tax
first, then takes the coupon off.

\smallskip
Same jacket, same coupon, same tax rate. Do the two of you pay the same
amount?

\writelines{2}
\end{hookbox}

\vspace{0.1in}

\boxguard[26]
\begin{notesbox}{1. Two Machines in a Row}
A \textbf{composition} runs two functions back to back: the output of the first
becomes the input of the second. We write it

\begin{center}
$(f \circ g)(x) = f\big(g(x)\big)$, \quad read ``$f$ of $g$ of $x$.''
\end{center}

The circle is \textit{not} multiplication. Here is the picture:

\smallskip
\begin{center}
\begin{tikzpicture}[x=1cm, y=1cm]
  \draw[thick, plum, fill=lilac, rounded corners=2pt]
    (-0.85,-0.55) rectangle (0.85,0.55);
  \node at (0,0) {\large $g$};
  \draw[thick, plum, fill=lilac, rounded corners=2pt]
    (4.15,-0.55) rectangle (5.85,0.55);
  \node at (5,0) {\large $f$};
  \draw[->, very thick, plum] (-2.6,0) -- (-0.95,0);
  \node[above, font=\small] at (-1.75,0.06) {$x$};
  \draw[->, very thick, plum] (0.95,0) -- (4.05,0);
  \node[above, font=\small] at (2.5,0.06) {$g(x)$};
  \draw[->, very thick, plum] (5.95,0) -- (8.2,0);
  \node[above, font=\small] at (7.05,0.06) {$f(g(x))$};
  \node[below, font=\small\itshape, text=slate] at (0,-0.62) {runs first};
  \node[below, font=\small\itshape, text=slate] at (5,-0.62) {runs second};
\end{tikzpicture}
\end{center}

\smallskip
Notice which letter runs first. In $f(g(x))$ the function written
\textit{last} is the one that goes \textit{first} --- exactly like innermost
parentheses in arithmetic.

\medskip
\begin{center}
\renewcommand{\arraystretch}{1.6}
\begin{tabularx}{0.94\linewidth}{@{} >{\bfseries}p{0.30\linewidth} X @{}}
\hline
The rule for today & Work from the \blank{2.2cm} out. The
  \blank{1.8cm} function runs first, and the number it produces is what you
  hand to the outer function. \\
\hline
$(f \circ g)(x)$ means & $f($ \blank{1.6cm} $)$, \quad \textit{not}
  $f(x) \cdot g(x)$. \\
\hline
\end{tabularx}
\end{center}
\end{notesbox}

\vspace{0.1in}

\boxguard[30]
\begin{notesbox}{2. Composite Values from a Table}
Everything a composition needs is a value to look up, so a table is enough ---
no formulas required.

\smallskip
\renewcommand{\arraystretch}{1.4}
\begin{center}
\begin{tabular}{|c|c|c|c|c|c|}
\hline
$x$ & 0 & 1 & 2 & 3 & 4 \\
\hline
$f(x)$ & 3 & 4 & 0 & 2 & 1 \\
\hline
$g(x)$ & 2 & 3 & 7 & 0 & 4 \\
\hline
\end{tabular}
\end{center}

\medskip
\textbf{Example 1.} Find $f(g(1))$ --- and \textit{write the middle number
down}.

\begin{enumerate}[label=\textbf{Step \arabic*.}, itemsep=3pt, leftmargin=4.2em]
  \item The inner function runs first: $g(1) = $ \blank{1.4cm}
  \item Hand that to the outer function: $f($ \blank{1.2cm} $) = $
        \blank{1.4cm}
\end{enumerate}

\medskip
\textbf{Example 2.} Now reverse the order: find $g(f(1))$.

\begin{enumerate}[label=\textbf{Step \arabic*.}, itemsep=3pt, leftmargin=4.2em]
  \item Now $f$ is on the inside: $f(1) = $ \blank{1.4cm}
  \item Hand that to $g$: $g($ \blank{1.2cm} $) = $ \blank{1.4cm}
\end{enumerate}

\smallskip
Same two functions, same input, \blank{2.4cm} answers. Composition is
\textbf{not commutative}: in general $f(g(x))$ \blank{1cm} $g(f(x))$.

\medskip
\textbf{Example 3 (a warning).} Try $f(g(2))$. First $g(2) = $ \blank{1.2cm}.
Now look for $f($ \blank{1.2cm} $)$ in the table --- it is
\blank{3cm}. The inner output has to be something the outer function
actually accepts; that is the \textbf{domain of a composite}.

\medskip
\textbf{Your turn.} Complete the bottom row.

\smallskip
\renewcommand{\arraystretch}{1.4}
\begin{center}
\begin{tabular}{|c|c|c|c|c|}
\hline
$x$ & 0 & 1 & 3 & 4 \\
\hline
$g(x)$ & 2 & 3 & 0 & 4 \\
\hline
$(f \circ g)(x) = f(g(x))$ & \blank{1.1cm} & \blank{1.1cm} & \blank{1.1cm}
  & \blank{1.1cm} \\
\hline
\end{tabular}
\end{center}
\end{notesbox}

\vspace{0.1in}

\boxguard[26]
\begin{notesbox}{3. Building a Composite Analytically}
With formulas instead of a table, the recipe is one sentence:

\begin{center}
To build $f(g(x))$, substitute $g(x)$ for \textbf{every} $x$ in $f(x)$.
\end{center}

You already did this in the warm-up, when $f(x) = 2x+1$ turned $f(x^2)$ into
$2x^2+1$. The only new part is the name.

\medskip
\textbf{Example 4.} Let $f(x) = 2x + 1$ and $g(x) = x^2 - 3$. Build both
composites.

\begin{work}
  f(g(x)) &= f(x^2 - 3) \\
          &= 2(x^2 - 3) + 1 \\
          &= 2x^2 - 6 + 1 \\
          &= 2x^2 - 5 \\
  g(f(x)) &= g(2x + 1) \\
          &= (2x + 1)^2 - 3 \\
          &= 4x^2 + 4x + 1 - 3 \\
          &= 4x^2 + 4x - 2
\end{work}

$(f \circ g)(x) = $ \blank{3.2cm} \qquad $(g \circ f)(x) = $ \blank{4.2cm}

\smallskip
These are not even the same shape of expression, which settles the question:
order \blank{2.4cm}.

\medskip
\textbf{Two routes, one number.} Check $f(g(2))$ both ways. Step by step:
$g(2) = 1$, so $f(1) = $ \blank{1.2cm}. From the formula:
$2(2)^2 - 5 = $ \blank{1.2cm}. A composite value can come from a table, a
graph, or an expression --- the answer does not care which.
\end{notesbox}

\vspace{0.1in}

\boxguard[26]
\begin{notesbox}{4. Order Matters --- the Checkout Counter}
Back to the jacket. Two functions act on the price $p$:

\begin{center}
\textbf{Coupon:} $c(p) = p - 10$ \qquad\qquad
\textbf{Tax:} $t(x) = 1.06x$
\end{center}

Neither one alone answers ``what do I pay?'' --- but chaining them does. That
is what composition is \textit{for}: linking two quantities no single formula
connects.

\medskip
\textbf{Example 5.} The jacket costs \$60. Compute both orders.

\begin{work}
  t(c(60)) &= t(60 - 10) \\
           &= 1.06(50) \\
           &= 53.00 \\
  c(t(60)) &= c(1.06 \cdot 60) \\
           &= 63.60 - 10 \\
           &= 53.60
\end{work}

Coupon first: \$\blank{1.6cm} \qquad Tax first: \$\blank{1.6cm}

\medskip
\textbf{Why?} Build both symbolically and the gap stops being a coincidence.

\begin{work}
  t(c(p)) &= 1.06(p - 10) \\
          &= 1.06p - 10.60 \\
  c(t(p)) &= 1.06p - 10
\end{work}

The two rules differ by \$\blank{1.4cm} at \textit{every} price, because
taxing first charges you tax on the \blank{2.6cm} you never actually paid.
\end{notesbox}

\vspace{0.1in}

\boxguard[26]
\begin{notesbox}{5. Decomposition, Identity, and Lesson 1.4}
\textbf{Decomposition} is composition run backwards: given one function, find
an inner and an outer piece that rebuild it.

\medskip
\textbf{Example 6.} Decompose $h(x) = (3x + 2)^4$. Ask: with a calculator,
what would you compute \textit{first}? That is the inner function.

\smallskip
Inner: $g(x) = $ \blank{2.6cm} \qquad Outer: $f(u) = $ \blank{2.2cm}

\smallskip
Check by substituting back: $f(g(x)) = $ \blank{3.4cm}

\medskip
\textbf{The identity function.} $I(x) = x$ hands back whatever you give it, so
composing with it changes nothing:

\smallskip
$I(f(x)) = $ \blank{1.8cm} \qquad and \qquad $f(I(x)) = $ \blank{1.8cm}

\smallskip
Hold on to this one --- next lesson is about finding a partner function that
turns $f$ into exactly $I$.

\medskip
\textbf{Lesson 1.4 was composition all along.} Every ``inside'' change you
made last lesson was a composition with a simple function.

\smallskip
\begin{center}
\renewcommand{\arraystretch}{1.6}
\begin{tabularx}{0.94\linewidth}{@{} >{\bfseries}p{0.26\linewidth} X @{}}
\hline
From Lesson 1.4 & Written as a composition \\
\hline
$g(x) = f(x+3)$ & $g(x) = f(h(x))$ \quad where $h(x) = $ \blank{2.4cm} \\
\hline
$g(x) = f(2x)$  & $g(x) = f(h(x))$ \quad where $h(x) = $ \blank{2.4cm} \\
\hline
\end{tabularx}
\end{center}

\smallskip
Now the reversal makes sense: the inner function gets to the $x$
\blank{2.6cm} $f$ does, so it is the \textbf{inputs} it rewrites.
\end{notesbox}

\vspace{0.1in}

\boxguard
\begin{practicebox}
\begin{enumerate}[itemsep=8pt]
  \item Use the table in section 2 to evaluate each composite.

        \begin{enumerate}[label=(\alph*), itemsep=4pt]
          \item $f(g(3)) = $ \blank{2cm}
          \item $g(f(3)) = $ \blank{2cm}
        \end{enumerate}

  \item Let $f(x) = x^2$ and $g(x) = x - 5$. Build both composites and
        simplify.

        \begin{work}
          f(g(x)) &= f(x - 5) \\
                  &= (x - 5)^2 \\
                  &= x^2 - 10x + 25 \\
          g(f(x)) &= g(x^2) \\
                  &= x^2 - 5
        \end{work}

        $(f \circ g)(x) = $ \blank{3.6cm} \qquad
        $(g \circ f)(x) = $ \blank{3cm}
\end{enumerate}
\end{practicebox}

\end{document}
